\documentclass[a4paper]{article}
% Español (México) con XeLaTeX: fontspec para fuentes Unicode, polyglossia
% para la división silábica y los nombres del documento en la variante mexicana.
\usepackage{fontspec}
\usepackage{polyglossia}
\setdefaultlanguage[variant=mexican]{spanish}
% Los nombres chinos van como «pinyin (original)»: xeCJK da a los caracteres
% Han su propia fuente; sin ella XeLaTeX los omite con un aviso.
% FandolSong no cubre algunos caracteres de nombres (U+73FA); WenQuanYi Zen Hei sí.
\usepackage[AutoFallBack=true]{xeCJK}
\setCJKmainfont{FandolSong-Regular.otf}
\setCJKfallbackfamilyfont{\CJKrmdefault}{WenQuanYi Zen Hei}
% polyglossia no traduce los nombres de listings.
\AtBeginDocument{\providecommand{\lstlistingname}{}\renewcommand{\lstlistingname}{Listado}%
  \providecommand{\lstlistlistingname}{}\renewcommand{\lstlistlistingname}{Índice de listados}}
\usepackage{amsmath, amssymb}
\usepackage{graphicx}
\usepackage[margin=2.5cm]{geometry}
\usepackage[most]{tcolorbox}
\usepackage{etoolbox}
\usepackage{listings}
\usepackage{hyperref}
\usepackage{booktabs}
\usepackage{subcaption}
\usepackage{float}

\newtcolorbox{knowledgebox}[1]{enhanced, colback=blue!5!white, colframe=blue!75!black, colbacktitle=blue!75!black, coltitle=white, fonttitle=\bfseries, title={#1}, attach boxed title to top left={yshift=-2mm, xshift=2mm}, boxrule=1pt, sharp corners}
\newtcolorbox{importantbox}[1]{enhanced, colback=yellow!10!white, colframe=yellow!80!black, colbacktitle=yellow!80!black, coltitle=black, fonttitle=\bfseries, title={#1}, sharp corners}
\newtcolorbox{warningbox}[1]{enhanced, colback=red!5!white, colframe=red!75!black, colbacktitle=red!75!black, coltitle=white, fonttitle=\bfseries, title={#1}, sharp corners}

\lstset{language=Python, basicstyle=\ttfamily\small, keywordstyle=\color{blue}, stringstyle=\color{red!60!black}, commentstyle=\color{green!60!black}, breaklines=true, frame=single, numbers=left, numberstyle=\tiny\color{gray}, captionpos=b, extendedchars=true}

\newcommand{\notetitle}{[LLM Agents SP25] LMs for Autoformalization \& Theorem Proving — Kaiyu Yang}
\newcommand{\noteauthors}{Elaboradas a partir de materiales públicos del curso}
\newcommand{\notedate}{2025}
\newcommand{\videochannel}{Berkeley RDI}
\newcommand{\videopublishdate}{2025}
\newcommand{\videoduration}{aproximadamente 60 minutos}
\newcommand{\videourl}{}
\newcommand{\videocoverpath}{cover.jpg}
\newcommand{\repourl}{https://github.com/hqhq1025/ai-course-notes}


% Footer with repo info
\usepackage{fancyhdr}
\pagestyle{fancy}
\fancyhf{}
\fancyfoot[C]{\thepage}
\fancyfoot[R]{\footnotesize\color{gray}\href{\repourl}{AI Course Notes}}
\renewcommand{\headrulewidth}{0pt}
\renewcommand{\footrulewidth}{0pt}

\begin{document}
\begin{titlepage}
\centering
{\Large Notas del curso\par}\vspace{1.2cm}
{\huge\bfseries \notetitle\par}\vspace{0.8cm}
{\large \noteauthors\par}\vspace{0.3cm}
{\large \notedate\par}\vspace{1.2cm}
\ifdefempty{\videocoverpath}{}{\includegraphics[width=0.82\textwidth,height=0.45\textheight,keepaspectratio]{\videocoverpath}\par}
\vfill
\begin{tcolorbox}[width=0.9\textwidth, colback=black!2!white, colframe=black!60, sharp corners]
\textbf{Autor o canal del video}: \videochannel\par
\textbf{Fecha de publicación}: \videopublishdate\par
\textbf{Duración del video}: \videoduration\par
\ifdefempty{\videourl}{}{\textbf{Enlace del video}: \href{\videourl}{\nolinkurl{\videourl}}\par}\par
\textbf{Página del proyecto}: \href{\repourl}{\nolinkurl{\repourl}}\par
\end{tcolorbox}
\end{titlepage}
\tableofcontents
\newpage

\section{Introducción: el razonamiento formal se encuentra con los modelos de lenguaje grandes}

Esta clase la impartió Kaiyu Yang, de Meta FAIR, y se centra en combinar el \textbf{razonamiento formal} (formal reasoning) con los modelos de lenguaje grandes, para impulsar la aplicación de la IA en las matemáticas y la verificación.

\begin{knowledgebox}{Por qué las matemáticas y la programación son indicadores centrales en la competencia entre LLM}
\begin{itemize}
    \item Las matemáticas y la programación son indicadores indirectos (proxy) del \textbf{razonamiento y la planeación complejos}
    \item La capacidad de razonamiento y planeación puede habilitar aplicaciones ilimitadas (planeación de viajes, gestión de agendas, etc.)
    \item Las matemáticas y la programación son \textbf{relativamente fáciles de evaluar}: en matemáticas se puede verificar la respuesta, y en código se pueden ejecutar pruebas unitarias
\end{itemize}
\end{knowledgebox}
\section{Dos técnicas principales para entrenar LLM de matemáticas}

\subsection{Fine-tuning supervisado (SFT)}

\begin{enumerate}
    \item El modelo base se preentrena con datos a escala de Internet
    \item Continúa el preentrenamiento con documentos matemáticos (StackOverflow, arXiv, etc.) $\to$ modelo matemático base
    \item Se hace fine-tuning con datos de preguntas y respuestas \textbf{que incluyen pasos detallados de solución} (puede incluir llamadas a herramientas, como Python)
\end{enumerate}

\begin{warningbox}{Los datos son el mayor cuello de botella}
La parte más costosa del pipeline de SFT es la recolección de datos: requiere anotación humana, limpieza y conversión de formato, y no puede automatizarse por completo.
\end{warningbox}

\subsection{Aprendizaje por refuerzo (RL)}

Cuando el conjunto de datos solo tiene la respuesta final y no los pasos intermedios, se puede entrenar con RL:
\begin{enumerate}
    \item El modelo genera una solución con pasos intermedios
    \item Se compara la respuesta final con la respuesta verdadera: recompensa 1 si es correcta, recompensa 0 si es incorrecta
    \item Se optimiza el modelo con algoritmos como GRPO para maximizar la recompensa
\end{enumerate}

\begin{importantbox}{La clave de RL: la verificabilidad}
RL depende de una señal de verificación confiable. Las respuestas numéricas se pueden comparar directamente, pero ¿cómo se verifica automáticamente la corrección de \textbf{un problema de demostración}? Aquí es exactamente donde reside el valor de los sistemas formales.
\end{importantbox}

\section{La brecha entre las matemáticas de competencia y las matemáticas de frontera}

\subsection{Limitaciones actuales de la capacidad matemática de los LLM}

\begin{itemize}
    \item O1/O3 tienen un desempeño sobresaliente en competencias preparatorias como AMC/AIME
    \item El benchmark Frontier Math contiene problemas de matemáticas de nivel de investigación; O3 resolvió más del 20\% (pero los problemas están diseñados para tener respuestas numéricas)
    \item Evaluación de Terence Tao: O1 todavía tiene dificultades en tareas de investigación matemática de frontera
\end{itemize}

\subsection{El problema de la generación de demostraciones}

\begin{warningbox}{El grave problema de los LLM al escribir demostraciones}
\begin{itemize}
    \item Competencia Putnam (nivel licenciatura): ChatGPT obtiene en su mayoría solo 1--2/10 puntos
    \item Evaluación USAMO: todos los modelos obtienen una puntuación de alrededor del 5\%
    \item Tipo de error: errores «simples» como la inversión del sentido de una desigualdad, ocultos dentro de demostraciones largas y sumamente difíciles de detectar
\end{itemize}
Causa raíz: en las matemáticas de frontera los datos son \textbf{escasos} y las demostraciones son \textbf{no verificables}.
\end{warningbox}

\section{Lean y la demostración formal de teoremas}

\subsection{Introducción a Lean}

Lean es un lenguaje de programación, un demostrador de teoremas y un asistente de demostración interactivo:
\begin{itemize}
    \item Define conceptos básicos como los números naturales y la suma
    \item Declara y demuestra teoremas (como la conmutatividad $a+b=b+a$)
    \item Una demostración se representa como un \textbf{árbol de demostración}: el nodo raíz es el teorema original, y cada paso lo descompone en subobjetivos más simples
    \item Los archivos se organizan en proyectos, y estos pueden reutilizarse entre sí (de forma similar a las dependencias de bibliotecas en el desarrollo de software)
\end{itemize}

\subsection{LeanDojo: infraestructura abierta para la demostración de teoremas}

\begin{importantbox}{LeanDojo (NeurIPS 2023)}
Proporciona conjuntos de datos, herramientas y modelos de código abierto:
\begin{itemize}
    \item Extrae alrededor de 10,000 teoremas y demostraciones a partir de código Lean escrito por humanos
    \item Extrae el estado de la demostración en cada paso, los lemas utilizados y sus definiciones
    \item Entrena \textbf{ReProver} (demostrador aumentado por recuperación): primero recupera los lemas relevantes y luego genera el siguiente paso
\end{itemize}
\end{importantbox}

\subsection{Mejora mediante Expert Iteration}

Se usa el demostrador actual para intentar demostrar nuevos teoremas $\to$ se recopilan las demostraciones exitosas $\to$ se agregan al conjunto de entrenamiento $\to$ se entrena un demostrador más fuerte $\to$ se itera, hasta que el desempeño se satura.

\subsection{El desafío fundamental de la demostración de teoremas}

\begin{warningbox}{Espacio de acciones infinito}
El go tiene un tablero finito de $19\times19$, pero el espacio de acciones de una demostración matemática es \textbf{prácticamente infinito}: es difícil cubrirlo solo con datos humanos, y la exploración con RL también es extremadamente difícil.
\end{warningbox}

\subsection{Caso: demostración de desigualdades (LiPS)}

En este dominio específico de las desigualdades de las olimpiadas de matemáticas:
\begin{itemize}
    \item Las acciones se dividen en \textbf{escalado} (aplicar lemas conocidos, finito y enumerable) y \textbf{reescritura} (transformar la fórmula, un espacio infinito que se deja al LLM)
    \item Un algoritmo simbólico enumera todos los escalados + el LLM se encarga de la reescritura + filtrado heurístico
    \item O1 Preview demuestra 0/20, O3 y DeepSeek-R1 demuestran cada uno 3--4/20, los medallistas de oro humanos 15/20, \textbf{LiPS demuestra 16/20}
    \item Se descubrió un método de demostración desconocido para los humanos (usando solo AM-GM junto con una reescritura ingeniosa)
\end{itemize}
\section{Formalización automática}

\subsection{Dos subproblemas}

\begin{enumerate}
    \item \textbf{Formalización de enunciados}: teorema no formal $\to$ enunciado formal
    \item \textbf{Formalización de demostraciones}: demostración no formal $\to$ demostración formal
\end{enumerate}

\subsection{El problema de evaluación de la formalización de enunciados}

Un mismo teorema puede tener varias formulaciones formales equivalentes (por ejemplo, «infinitos números primos» tiene varias escrituras equivalentes), pero comprobar la equivalencia lógica entre dos enunciados formales \textbf{es indecidible en el caso general}.

\subsection{La brecha de razonamiento en la formalización de demostraciones}

Las demostraciones humanas, aun siendo rigurosas, a menudo tienen partes «dejadas al lector»; el modelo debe llenar por sí mismo estas \textbf{brechas de razonamiento}.

\subsection{Caso de estudio: formalización automática en la geometría euclidiana}

Se elige la geometría euclidiana como dominio restringido, aprovechando las particularidades del dominio para volver resoluble un problema originalmente irresoluble.

\subsection{Resumen de la sección}

El reto central de la demostración de teoremas es cómo explorar de manera eficiente un espacio de acciones infinito. Las intuiciones específicas del dominio pueden estructurar el espacio de acciones, pero la generalización sigue siendo un problema abierto.
\section{Síntesis y ampliación}

\subsection{Mensaje central}

\begin{enumerate}
    \item La capacidad matemática actual de los LLM depende principalmente del SFT y el RL, limitada por el volumen de datos y la posibilidad de verificación
    \item El razonamiento formalizado (Lean) ofrece una señal de verificación perfecta, lo que puede aliviar la escasez de datos y la dificultad de evaluación
    \item LeanDojo ofrece una infraestructura de código abierto para la demostración de teoremas
    \item Los métodos específicos de dominio (como LiPS) pueden superar a los humanos y a los LLM de vanguardia en tareas específicas
    \item La formalización automática (de enunciados y de demostraciones) presenta retos particulares en cada caso
\end{enumerate}

\subsection{Lecturas adicionales}
\begin{itemize}
    \item LeanDojo (NeurIPS 2023): datos y herramientas de código abierto para la demostración de teoremas
    \item LiPS: método híbrido simbólico y neuronal para la demostración de desigualdades
    \item Artículo de postura ``AI for Math'' (arXiv)
    \item Referencia Frontier Math (Epoch AI)
    \item DeepSeek-R1: entrenamiento con RL de un modelo de razonamiento
\end{itemize}

\end{document}
